\chapter{Introduction}

\section{Background}
The standard joints of machine design, such as the knuckle joint, the cotter joint and the shaft coupling, are not designed by judgement: each is sized by a fixed sequence of formulas and verified against a fixed list of failure modes, as in Bhandari. The sequence is repetitive, since a change in load invalidates every dimension. Computer aided design software removes the labour of drawing but not of calculation; a parametric modeller does not object if a pin is too slender for the bending moment upon it.

\section{Motivation}
The motivation is the gap between stating a requirement and obtaining a verified design. At present a designer must translate a requirement such as \emph{a knuckle joint to carry 50 kN} into formulas by hand, check several failure modes, iterate wherever a check fails, and only then begin modelling and drafting.

A system that accepts the requirement as written and returns a verified design, with solid models and a drawing sheet, would remove the repetition without removing the engineering, provided the calculation itself stays checkable.

\section{Objectives}
The objectives of the present work are as follows.
\begin{enumerate}
  \item To develop a system that accepts a design requirement in ordinary English and extracts the joint type, the applied load and the allowable stresses.
  \item To size the joint using established design formulas and verify every applicable failure mode against the allowable stress, iterating wherever a check is not satisfied.
  \item To generate exact solid models of every component and of the assembly in a neutral exchange format that opens in commercial CAD software.
  \item To generate a first angle projection drawing sheet carrying dimensions and a title block in conformity with Indian drawing practice.
  \item To verify every proportion and failure check used by the system against published design literature, recording any disagreement with a source explicitly.
  \item To make the system usable through a web browser without any installation.
\end{enumerate}

\section{Scope of the Work}
The work covers six standard joints under static loading. Fatigue, stress concentration factors, tolerances and fits, and manufacturing processes are outside its scope; the finite element facility is a single linear static cross check, not a complete analysis.

\section{Organisation of the Report}
Chapters 2 to 4 review the literature, state the problem and describe the case studies; Chapters 5 and 6 give the methodology and its implementation; Chapter 7 presents the results and their verification, and Chapter 8 concludes.
